% section: 解法の決定

\section{解法の決定}\label{節：解法の決定}
  以下にここまでで扱った解法を体系的にまとめあげ,実際に問題と相対した時にどこに注目すればよいか,なにを考えればよいかをまとめた.
  \chapterref{章：漸化式の解説}で目指してきたように,式の特徴を手掛かりとして,既知の「きれい」な形に帰着させることを考えよう.

  チャートは「式全体」「数列の項」「足されている数」の3つの観点から,着目点と試せる変形を整理したものである.
  すべてを順番に確認する必要はなく,式の中で目立つ特徴から必要な項目を使えばよい.
  一つの式に複数の特徴が当てはまることもある.どの特徴を使い,どの順番で変形するかは,この後の問題演習で練習していこう.
  変形した後にも,新しい式について同じ着目点を確認するとよい.

  \noindent\textbf{【観察から変形まで】}\par
  次の4段階を,必要に応じて繰り返す.チャートの出口は候補であり,一つの正解に固定するものではない.
  \begin{enumerate}
    \item \textbf{観察}：総和,分母,倍率,付加項のうち,何が目立つかを具体的な式で指す.
    \item \textbf{候補}：その特徴が変形後に消えるか,一定になるかを考える.複数の候補があってよい.
    \item \textbf{条件}：割る量の非零,対数の真数の正値,初項による定数解の分岐を確かめる.
    \item \textbf{変形と評価}：新しい量の更新を一行計算し,何が簡単になったかを見る.新しい式を観察し直し,必要なら別の候補へ戻る.
  \end{enumerate}
  候補が正しくても計算が長くなる場合はある.条件の見やすさと計算量も比較しよう.
  一般項が得られたら,初期条件と元の更新式を確認する.

  \begingroup
  \clubpenalty=10000
  \widowpenalty=10000
  \medskip
  \noindent \textbf{【チャート1：式全体について】}
  \begin{enumerate}
    \item \textbf{総和$S_n$があるか}

      $S_n=\sum_{k=1}^{n}a_k$なら,\,$S_{n+1}-S_n=a_{n+1}$である.
      添字を一つずらした式との差を取り,総和を消去できないか考える（総和型）.

    \item \textbf{数列の項を含む式が分母にあるか}

      逆数を取って簡単にならないか考える.
      定数係数の1次式同士の商なら,定数を引いて逆数型へ帰着させる方法もある（分数型）.
      分母が$0$にならないことを確認する.

    \item \textbf{何項間か.添字は隣接しているか}

      例えば$a_{n+2}$と$a_n$の関係なら,奇数番目と偶数番目を別々の数列として考えられる（非隣接2項間漸化式型）.
      また,\,$a_{n+2}=qa_n$は
      \[
        a_{n+2}=0\cdot a_{n+1}+qa_n
      \]
      と書ける.定数係数で,各項が1次の和として現れるなら,欠けた項を係数$0$で補い,隣接3項間・4項間の解法を検討する.

    \item \textbf{展開・因数分解・通分で簡単になるか}

      見た目が複雑でも,整理すると基本的な形が現れることがある.
      分数を約分する際には,消す因子が$0$になる場合を別に確認する.

    \item \textbf{同じ式のまとまりが繰り返し現れるか}

      例えば$a_n+n$と$a_{n+1}+(n+1)$が現れるなら,\,$b_n=a_n+n$と置いて関係を調べる.

    \item \textbf{既知の公式に似ているか}

      $a_n^2-2$のような形なら,積が$1$となる二つの数の和を使う置換を考える（基本対称式型）.
      $2a_n^2-1$や$\displaystyle \frac{2a_n}{1-a_n^2}$は,三角関数の倍角公式を思い出す手掛かりになる（三角関数型）.
      $\displaystyle \frac12\left(a_n+\frac{c}{a_n}\right)$のような平均の形なら,ニュートン法型を検討する.
      見た目に加えて,初項や置換の条件も確認する.

    \begin{samepage}
    \item \textbf{複数の数列の関係が与えられているか}\\*[3pt]
      定数係数の線形連立漸化式なら,数列の和・差などを使って一つずつの関係に分けられないか考える（連立型）.
    \par
    \end{samepage}

    \item \textbf{整数部分を表す記号があるか}

      まず,整数部分の記号が項の値・添字・付加項のどこに作用しているかを見る（ガウス型）.
      項の値なら整数性や余り,添字なら偶奇や群への分割,付加項なら同じ値を取る範囲が手掛かりになる.
      床関数を外せるかを調べるときは,整数性と上下の不等式を確かめる.
  \end{enumerate}

  \medskip
  \noindent \textbf{【チャート2：数列の項について】}

  係数については,主に$a_{n+1}=p_na_n+f(n)$の形に整理できる場合を考える.
  \begin{enumerate}
    \item \textbf{係数は定数か,\,$n$に依存するか}

      \begin{itemize}[beginpenalty=10000]
        \item 定数の場合は,等比型や特性方程式型などを考える.係数が$1$なら,階差に注目する.
        \item $n$に依存する場合は,その係数の積が求めやすいか調べる（階比型）.
              足されている数もある場合は,適切な数列で割って係数をそろえられないか考える.
              例えば係数が$\displaystyle \frac{n+1}{n}$なら,\,$b_n=\displaystyle \frac{a_n}{n}$という置換が候補になる.
      \end{itemize}
      係数だけで判断せず,足されている数についてもチャート3を確認する.

    \item \textbf{累乗の指数はどうなっているか}

      \begin{itemize}[beginpenalty=10000]
        \item すべて1次か.2次以上なら,因数分解や$a_n^2$などのまとまりを使う置換を考える.
        \item 指数は定数か,\,$n$に依存するか.対数を取ると,指数を係数として取り出せる場合がある（対数型）.
      \end{itemize}
      実数の対数を取る各量は,正であることを確認する.

    \item \textbf{数列の項同士の積や商があるか}

      $a_na_{n+1}$などの積があるなら,因数分解や置換でまとまりを作れないか考える.
      積・商・累乗を対数で和・差・定数倍に変えられるかも調べる.
      ただし,和の対数を各項の対数の和に分けることはできない.
  \end{enumerate}

  \medskip
  \noindent \textbf{【チャート3：足されている数について】}

  ここでは,数列の項を含まずに足されている部分を「付加項」と呼ぶ.
  以下では$a_{n+1}=p_na_n+f(n)$の形を考え,\,$f(n)$を調べよう.
  \begin{enumerate}
    \item \textbf{付加項はないか,定数か}

      付加項がなければ,係数の積から求める方法を考える（等比型・階比型）.
      付加項も係数も定数なら,定数を引いて消せないか考える（特性方程式型）.
      ただし$a_{n+1}=a_n+d$なら,そのまま等差型として解ける.

    \item \textbf{$n$に依存するなら,どのような形か}

      \begin{itemize}[beginpenalty=10000]
        \item \textbf{指数の形}：$cq^n$などなら,\,$q\ne0$を確認し,\,$q^n$などで割って指数の部分を取り除くことを考える.
        \item \textbf{$n$の多項式}：$n^2+2n$などで係数が定数なら,多項式を引いて付加項を消す方法を考える（階差等比複合型）.
        \item \textbf{多項式と指数の積}：$n2^n$などなら,まず指数の部分で割り,残った多項式をどう扱うか考える.
      \end{itemize}
      例えば$n+2^n$のように複数の形が混ざっていれば,それぞれの部分を見る.

    \item \textbf{総和を求めやすい形にできるか}

      階差型へ帰着できたら,付加項の総和を考える.
      $n(n+1)$のような連続整数の積や,隣り合う項の差として表せないか調べる.
      差の形にできれば,総和を取ったときに途中の項が消える.
  \end{enumerate}

  \medskip
  \noindent \textbf{【補足：初項も手掛かりになる】}

  例えば初項が$\displaystyle \frac{1}{\sqrt{6}-\sqrt{2}}$だとしよう.
  わざわざこのような数字が設定されたということはきっと裏がある.
  実際これは$\cos 15^\circ$または$\sin 75^\circ$の値である.
  そこから考えて,この漸化式は三角関数型のものではないか,と予想するのである.
  ただし,初項だけで解法が決まるわけではない.漸化式の形と合わせて判断しよう.
  \par
  \endgroup

% 解法の選択を、計算の前に確認する短い練習。
\subsection{操作を選ぶ練習}\label{節：操作を選ぶ練習}
  ここでは\bookterm{general}{一般項}まで一気に計算する前に,新しい量を選び,その更新だけを確かめる.
  引く・割る・\bookterm{logarithm}{対数}を取るという操作が,それぞれ何を簡単にするかを比較しよう.

  \noindent\bookblocklink{operation-choice}{problem}{answer}{【問題】}\par\nobreak
  \begin{enumerate}[format=\bookitemlink{operation-choice}{problem}{answer}]
    \item $a_1=1$,\,$a_{n+1}=3a_n+2n+1$とする.
          定数を引くだけで付加項を消せるか.
          $g(n)=un+v$を引く場合の$u,v$を求め,\,$b_n=a_n-g(n)$の更新式と初項を示せ.
    \item $a_1=2$,\,$a_{n+1}=\frac{n+2}{n}a_n+(n+1)(n+2)$とする.
          $b_n=\frac{a_n}{n}$と$c_n=\frac{a_n}{n(n+1)}$を候補とし,それぞれの更新式を求めよ.
          どちらで$a_n$の係数を$1$にそろえられるか.
    \item $a_1=-2$,\,$a_{n+1}=2^{2n+1}a_n$とする.
          $b_n=\log_2a_n$は使えるか.
          \bookterm{logarithm}{対数}を使わず,\,$c_n=\frac{a_n}{2^{n^2}}$の更新式と初項を求めよ.
  \end{enumerate}

  第2問は,\sectionref*{節：階比型}で倍率をそろえた操作を,付加項がある場合に使う練習である.
  一般に
  \[
    a_{n+1}=p_na_n+q_n
  \]
  に対して,既知の非零の\bookterm{sequence}{数列}$h_n$が$h_{n+1}=p_nh_n$を満たすなら,
  \[
    b_n=\frac{a_n}{h_n},\qquad
    b_{n+1}-b_n=\frac{q_n}{h_{n+1}}
  \]
  となる.倍率をそろえる操作の後に,\bookterm{difference}{階差}を足し戻す操作が続く.
  $p_n\ne0$なら,\,$h_1=1$,\,$h_n=\prod_{k=1}^{n-1}p_k$を候補にできる（空積は$1$）.
  係数が零になる場合はこの割り算を使わず,その段階の更新を別に確認する.
  また,残った和がいつでも簡単に計算できるとは限らない.

  \noindent\bookblocklink{operation-choice}{answer}{problem}{【解答と選択の理由】}\par\nobreak
  \begin{enumerate}[format=\bookitemlink{operation-choice}{answer}{problem}]
    \item $a_n-\alpha$では付加項の$2n$が残る.
          $g(n+1)=3g(n)+2n+1$から
          \[
            u=3u+2,\qquad u+v=3v+1
          \]
          となり,\,$u=v=-1$である.
          従って$b_n=a_n+n+1$により$b_{n+1}=3b_n$,\,$b_1=3$を得る.
          引く量を定数から一次式へ変えたことで,付加項を消せた.

    \item 候補ごとに一歩だけ更新すると,
          \[
            b_{n+1}=\frac{n+2}{n+1}b_n+n+2
          \]
          と
          \[
            c_{n+1}=c_n+1,\qquad c_1=1
          \]
          を得る.
          $h_n=n(n+1)$なら
          \[
            \frac{h_{n+1}}{h_n}=\frac{n+2}{n}
          \]
          が元の倍率と一致する.
          この一致を手掛かりに割る因子を選ぶ.
          $n\ge1$なので分母は非零であり,\,$c_n=n$から$a_n=n^2(n+1)$も得られる.

    \item 正の倍率で更新するので,各項は負であり,\,$\log_2a_n$は実数として定義できない.
          一方,\,$2n+1=(n+1)^2-n^2$から
          \[
            c_{n+1}=c_n,\qquad c_1=-1
          \]
          となる.従って$a_n=-2^{n^2}$を得る.
          \bookterm{normalize}{正規化}で割る$2^{n^2}$は常に非零であり,元の項が正である必要はない.
          初項が正の比較例は\bookproblemref{practice-basic}{42}{標準問42}へ.
  \end{enumerate}

